\chapter{Introdu\c{c}\~ao}
\label{pre:introducao}

Nos \'ultimos anos, as redes neurais de grafos tornaram-se uma importante \'area de pesquisa em computa\c{c}\~ao e ci\^encia de dados, ganhando destaque tanto na academia quanto na ind\'ustria \cite{kipf2017}. Esta abordagem destaca-se por sua capacidade de processar grandes quantidades de dados interligados, fornecer \textit{insights} valiosos e identificar padr\~oes complexos em sistemas tamb\'em complexos. Suas aplica\c{c}\~oes foram encontradas em uma ampla gama de disciplinas, incluindo an\'alise de dados educacionais, em que o reconhecimento de padr\~oes desempenha papel relevante na identifica\c{c}\~ao de tend\^encias e na melhoria da experi\^encia acad\^emica \cite{wang2021}.

O reconhecimento de padr\~oes nos dados educacionais \'e fundamental, pois pode contribuir para prevenir a evas\~ao escolar e melhorar o progresso acad\^emico dos alunos. No entanto, as limita\c{c}\~oes de m\'etodos tradicionais de reconhecimento de padr\~oes frequentemente dificultam a detec\c{c}\~ao precoce de problemas e a implementa\c{c}\~ao de interven\c{c}\~oes eficazes \cite{jain2020}. Nesse contexto, as redes neurais de grafos surgem como uma ferramenta promissora para melhorar o reconhecimento de padr\~oes em dados educacionais.

Sua capacidade de modelar rela\c{c}\~oes complexas entre diferentes vari\'aveis e entidades pode fornecer m\'etodos de reconhecimento mais precisos e eficazes, permitindo a detec\c{c}\~ao precoce de problemas acad\^emicos e a aplica\c{c}\~ao de interven\c{c}\~oes mais direcionadas.
